\chapter{Abstract}

\noindent Driving a two-armed robot from a distance, through a single hand-held haptic interface, is difficult for reasons that are structural rather than incidental. The operator perceives the workcell only through the robot, commands one arm at a time through a device whose workspace is a fraction of the robot's, and must keep both arms clear of the environment and of each other while doing so. Assistance can relieve part of that load: a layer that infers which object the operator is reaching for can steer the arm toward it and can push the operator's hand toward it through the device. The difficulty is that assistance which is confidently wrong is worse than no assistance at all, and assistance that overrides the operator removes exactly the judgement the operator was present to supply.

This thesis addresses that trade-off by separating what the robot can decide from what only the operator can. A reactive layer computes, at every instant, the best locally admissible motion toward each candidate goal, and enforces collision avoidance and joint limits as hard constraints of a nominal barrier condition on whatever command finally reaches the arms, so that no assistance and no operator input can produce an unsafe motion. That layer is deliberately local, and its locality is what defines the division of labour: it cannot know which of several objects is intended, which side of an obstacle to pass, or when a temporary detour away from the goal is the only way out of a dead end. Those decisions remain with the operator. Assistance is therefore arbitrated on agreement --- it contributes while the operator is moving with it, and withdraws as soon as the operator moves against it --- so the robot supplies precision and safety while the operator supplies intent and judgement.

The architecture was implemented on a bimanual mobile manipulator and validated end to end on the physical robot, from camera-perceived obstacle geometry through teleoperated grasping under barrier-enforced collision avoidance. The evaluation crossing the two classical teleoperation mappings --- position control with indexing and spring-centred rate control --- with three assistance strategies --- haptic guidance, autonomous blending of the operator's command with the inferred policy, and both together --- over two tasks that differ in whether they force the two arms to cooperate, was run with human participants in a matched simulation environment.

Across the factorial comparison, the assistance strategy, not the motion mapping, is what moves objective performance: reference blending is what shortens the task and steadies the hand, force guidance alone is consistently the weakest condition in both mappings, and no significant difference was found between blending alone and blending with guidance. The minimum clearance recorded shows no significant difference across conditions, consistent with the barrier, rather than the operator or the assistance, setting that floor. Participants agree with that objective picture only in part: they rate the blended conditions as more controllable and more trustworthy than guidance alone, even though blending is exactly the channel that takes some authority away from the hand, while the choice of motion mapping, despite separating the two objectively, is something participants report noticing hardly at all. On the physical robot, the same controller and arbitration carried a real teleoperated grasp to completion and, on a second attempt, recorded the architecture's one safety violation, when a clearance-blind retreat re-armed a deliberately suspended collision barrier onto an unresolved overlap before the operator stopped the robot; that single case does not settle how the architecture behaves at hardware scale, and is treated here as a limitation rather than as a conclusion.
